\section{边郡财政、迁徙与军人社会：东汉后期的积累}
\label{sec:eastern-han-border-finance-society}

东汉的西、北边防并不是宫廷危机之外的背景。永初元年（一〇七），西羌在凉州、并州一带再起，朝廷必须长期投入军费并安置流民、降附者和迁徙者。军功叙事较容易写出一支军队何时出发，却没有留下同等连续的粮秣、逃户和家庭迁移账簿。较为清楚的是，战乱使凉州成为持续的财政负担，也使地方社会中军人、郡县官、被征发者与部众之间的关系不断重组。

\subsection{以围剿维持道路的代价}

一六七年至一六九年，段颎在陇右连续进军，东羌多有击散或降附，战事伴随严酷杀戮。短期恢复部分通道与郡县控制，不等于冲突的条件已经消失：征发、牧地、地方联盟和军费仍会影响下一次动员。斩获、迁徙和“平定”的范围在传世叙事中尤其需要克制；把“羌乱”当作单一敌人，也会抹去军民、地方首领和不同部众的相异处境。

一七七年檀石槐领导下的鲜卑多次侵扰北边，汉军出击受挫。北方威胁从匈奴转向多部鲜卑联盟之际，边郡已在羌乱和财政紧张中难以维持机动防御。史书的“鲜卑”同样是历史分类而非稳定的现代族群标签；部众的统合程度、兵数与损失均不能由概称推定。西部与北部同时需要募兵、粮运和地方合作，才使边防的压力逐渐传导到中枢的人事与财政。

\subsection{凉州、幽州与地方军权}

一八四年凉州军民与羌胡部众起事，北宫伯玉、边章、韩遂等先后成为叛军领袖。它发生在长期边防征发、羌乱遗绪和军人利益冲突交织之处；黄巾危机又使朝廷难以将资源集中于西北。此后中央委任董卓等人带兵镇压，西北逐渐形成能够自筹、招募并与朝廷谈判的军事集团。地方军权并非一八四年才突然出现，而是在边郡长久的征发和安置中获得了制度化的用途。

一八六年张纯、张举在幽州起兵，东北边郡与乌桓的关系再度动荡。中央忙于黄巾和凉州时，幽州地方军人与边郡网络可以利用征发压力发动反叛，公孙瓒等将领的地位随之上升。凉州、幽州的情形不能合并为同一种“割据”：道路、牧地、邻近部众和既有军府不同；但两地都显示，地方一旦掌有兵员、给养和对外关系，洛阳的任命并不能自动恢复直接控制。

\subsection{许都屯田与北方侧翼}

一九六年曹操迎献帝至许县，随后以朝廷名义发诏并组织屯田。迎驾使曹氏取得皇帝名义，屯田则将粮仓、军队、行政人员和可征发的人户编进较持续的军政秩序。屯田初创的确切时间、规模和曹操个人作用仍有争议，不宜把它写成一项从上至下无阻施行的制度；它至少说明，竞争者要利用皇帝名义，仍须解决军人和居民的日常供给。

二〇七年曹操北征乌桓，在白狼山击破蹋顿，袁氏残部失去北方依托。胜负与北征的发生可靠，行军路线、参战部众及战后安置却不能简化为“完全臣服”。幽州通道被纳入曹氏控制，既为此后南下减少侧翼风险，也把边疆人群的安置和地方军事关系留给新的北方政权处理。东汉终局由此不是单一宫廷崩溃，而是边郡财政、迁徙、军事集团与皇帝名义被不同力量重新组合。

\subsection{来源与空白}

《后汉书》帝纪、〈西羌传〉、〈鲜卑传〉、董卓传和乌桓鲜卑列传构成本节的年序主干，《后汉纪》可校读异文；许都屯田还须同《三国志》武帝纪、任峻传相互检验。它们集中保存将领、诏令和战事，难以连续记录被迁徙者、军户和普通边民。本节因此把军费、道路、安置和招募作为可见的制度压力，同时保留部众构成、杀伤数字和地方执行程度的材料边界。
